\bta{2017}


\begin{enumerate}
	%\renewcommand{\labelenumi}{\arabic{enumi}.}
	% A(\Alph) a(\alph) I(\Roman) i(\roman) 1(\arabic)
	%设定全局标号series=example	%引用全局变量resume=example
	%[topsep=-0.3em,parsep=-0.3em,itemsep=-0.3em,partopsep=-0.3em]
	%可使用leftmargin调整列表环境左边的空白长度 [leftmargin=0em]
	\item
	(共30分)一个粒子在一维无限深势阱
	$$
	V(x)=\left\{\begin{array}{cc}
		0 & 0<x<a \\
		\infty & x<0, x>a
	\end{array}\right.
	$$
	中运动,其势阱内定态波函数为$\psi_{n}(x)=\sqrt{\frac{2}{a}} \sin \frac{n \pi x}{a}$。
	设$ t=0 $时刻粒子势阱内的波
	函数为$\Psi(x, 0)=A x(a-x)$
	\begin{enumerate}
		%\renewcommand{\labelenumi}{\arabic{enumi}.}
		% A(\Alph) a(\alph) I(\Roman) i(\roman) 1(\arabic)
		%设定全局标号series=example	%引用全局变量resume=example
		%[topsep=-0.3em,parsep=-0.3em,itemsep=-0.3em,partopsep=-0.3em]
		%可使用leftmargin调整列表环境左边的空白长度 [leftmargin=0em]
		\item
		求归一化常数$ A $。

		\item 
		求$ t>0 $时刻粒子的波函数$ \Psi(x,t) $。

		\item 
		求$ t>0 $时刻粒子处于系统基态及第一激发态的概率。

		\item 
		计算$ t>0 $时刻粒子位置的平均值。（提示：$\sum\limits_{\mathstrut m=1}^{\mathstrut \infty} \dfrac{1}{(2 m-1)^{6}}=\dfrac{\pi^{6}}{960}$）
	
		
		
		
		
		
		
			\end{enumerate}

	\banswer{
		\begin{enumerate}
			\item 归一化条件给出
			\[
				1=A^2\int_0^a x^2(a-x)^2dx=A^2\frac{a^5}{30},
			\]
			故可取
			\[
				\boxed{A=\frac{\sqrt{30}}{a^{5/2}}}.
			\]

			\item 展开系数为
			\begin{align*}
				c_n&=\int_0^a\psi_n^*(x)\Psi(x,0)dx\\
				&=\begin{cases}
				\dfrac{8\sqrt{15}}{n^3\pi^3},&n=1,3,5,\ldots,\\[5pt]
				0,&n=2,4,6,\ldots.
				\end{cases}
			\end{align*}
			因此
			\[
				\boxed{\Psi(x,t)=\sum_{\substack{n=1\\ n\ {\rm odd}}}^{\infty}
				\frac{8\sqrt{15}}{n^3\pi^3}
				\sqrt{\frac2a}\sin\frac{n\pi x}{a}\,
				e^{-iE_nt/\hbar}},
				\qquad E_n=\frac{n^2\pi^2\hbar^2}{2ma^2}.
			\]
			提示中的级数恰好验证
			$\sum_{n\ {\rm odd}}|c_n|^2
			=(960/\pi^6)\sum_{n\ {\rm odd}}n^{-6}=1$。

			\item 基态概率和第一激发态概率为
			\[
				\boxed{P_1=|c_1|^2=\frac{960}{\pi^6}},
				\qquad
				\boxed{P_2=|c_2|^2=0}.
			\]
			它们均不随时间变化。

			\item 初态关于势阱中心 $x=a/2$ 对称；展开中只有关于该中心为偶宇称的奇数编号本征态，时间演化不破坏此宇称。因此任意时刻的概率密度均满足
			$|\Psi(a-x,t)|^2=|\Psi(x,t)|^2$，从而
			\[
				\boxed{\langle x\rangle_t=\frac a2}.
			\]
		\end{enumerate}
	}


	\item
(共$ 30 $分)一个能量为$ E $的粒子,沿$ x $轴从左侧入射。阶梯势垒为
$$
V(x)=\left\{\begin{array}{ll}
	0, & x<0 \\
	V_{a}, & 0<x<a \\
	V_{b}, & x>a
\end{array}\right.
$$
如图所示。
\begin{figure}[h!]
	\centering
	\includesvg[width=0.33\linewidth]{picture/svg/811-001}
\end{figure}
\begin{enumerate}
	%\renewcommand{\labelenumi}{\arabic{enumi}.}
	% A(\Alph) a(\alph) I(\Roman) i(\roman) 1(\arabic)
	%设定全局标号series=example	%引用全局变量resume=example
	%[topsep=-0.3em,parsep=-0.3em,itemsep=-0.3em,partopsep=-0.3em]
	%可使用leftmargin调整列表环境左边的空白长度 [leftmargin=0em]
	\item
	取入射波函数为$\psi_{i}(x)=e^{i k_{1} x}$,  $k_{1}=\sqrt{\frac{2 m E}{\hbar^{2}}}$,求入射粒子流密度。
	
	\item 
	选择题:若$ E=V_{b} $,透射系数$ T $是多少?
\threechoices
{$1$}
{$0<T<1$}
{$0 $}

\item 
选择题:要增大透射系数$ T $,以下哪种做法是正确的?


\threechoices
{$ V_{b} =E $,减小$ V_{a} $,但保持$ V_{b}<V_{a} $}
{减小$ V_{b}( V_{b} <E) $,但保持$ V_{a} $不变}
{$ V_{b} =E $,减小$ V_{a} $,使$ V_{a}< V_{b} $}


\item 
选择题:设$ E>V $,反射波函数为$\psi_{r}(x)=B e^{-i k_{1} x}$ ,透射波函数为$\psi_{t}(x)=C e^{i k_{2} x}$，其中$k_{1}=\sqrt{\frac{2 m E}{\hbar^{2}}},  k_{2}=\sqrt{\frac{2 m\left(E-V_{b}\right)}{\hbar^{2}}}$。问透射系数$ T $是多少?

\threechoices
{$ |C|^{2} $}
{$1-|B|^{2}$}
{以上都不是}





	
	
	
	
	
	
	
\end{enumerate}

	\banswer{
		\begin{enumerate}
			\item 一维概率流密度为
			\[
				j=\frac{\hbar}{2mi}\left(\psi^*\frac{d\psi}{dx}
				-\psi\frac{d\psi^*}{dx}\right).
			\]
			对 $\psi_i=e^{ik_1x}$，
			\[
				\boxed{j_i=\frac{\hbar k_1}{m}
				=\sqrt{\frac{2E}{m}}}.
			\]

			\item 当 $E=V_b$ 时，右侧区域的波数
			$k_2=\sqrt{2m(E-V_b)}/\hbar=0$，透射概率流为零，故
			\[
				\boxed{T=0},
			\]
			应选第三项。

			\item 在其余参数不变时，把 $V_b$ 从 $E$ 降到 $V_b<E$，可使右侧出现传播波并使 $T$ 从零变为非零。因此应选第二项。

			\item
			\[
				T=\frac{j_t}{j_i}
				=\frac{k_2}{k_1}|C|^2.
			\]
			它一般不等于 $|C|^2$；对实势场又由概率流守恒得到
			$T=1-|B|^2$。故应选第二项。
		\end{enumerate}
	}


\item 
(共30分)
\begin{enumerate}
	%\renewcommand{\labelenumi}{\arabic{enumi}.}
	% A(\Alph) a(\alph) I(\Roman) i(\roman) 1(\arabic)
	%设定全局标号series=example	%引用全局变量resume=example
	%[topsep=-0.3em,parsep=-0.3em,itemsep=-0.3em,partopsep=-0.3em]
	%可使用leftmargin调整列表环境左边的空白长度 [leftmargin=0em]
	\item
设质量为$ m $的粒子在一维势场$ V(x) $运动,试证明  Ehrenfest 定理:
$$\left\{\begin{aligned}
	\frac{d}{d t}\langle x\rangle&=\frac{\langle p\rangle}{m} \\ \frac{d}{d t}\langle p\rangle&=-\left\langle\frac{\partial V}{\partial x}\right\rangle
\end{aligned}\right.$$
其中，$\langle F\rangle=\int \psi^{*}(x, t) F(x, p) \psi(x, t) d x$表示算符$ F $的平均值。

\item 
设二维各向异性谐振子在均匀电场$\left(\varepsilon_{x}, \varepsilon_{y}\right)$中运动,哈密顿量为
$$
H=\frac{p_{r}^{2}}{2 m}+\frac{p_{y}^{2}}{2 m}+\frac{1}{2} m \omega_{x}^{2} x^{2}+\frac{1}{2} m \omega_{y}^{2} y^{2}-q\left(\mathcal{E}_{x} x+\mathcal{E}_{y} y\right)
$$
$ 2 \ m^{2}m2 $
其中$ q $为粒子电荷。试写出无电场情形下体系的束缚态能级,并以$ q $为参数用
Hellmann-Feynman 定理计算存在电场情形下体系的束缚态能级。




	
\end{enumerate}

\banswer{
	\begin{enumerate}
		\item 对任意不显含时间的算符 $F$，
		\[
			\frac{d\langle F\rangle}{dt}
			=\frac{i}{\hbar}\langle[H,F]\rangle.
		\]
		取 $H=p^2/(2m)+V(x)$。由
		$[H,x]=-(i\hbar/m)p$ 与 $[H,p]=i\hbar\,\partial V/\partial x$，立即得到
		\[
			\boxed{\frac{d\langle x\rangle}{dt}=\frac{\langle p\rangle}{m}},
			\qquad
			\boxed{\frac{d\langle p\rangle}{dt}
			=-\left\langle\frac{\partial V}{\partial x}\right\rangle}.
		\]

		\item 无电场时能级为
		\[
			E_{n_xn_y}(0)=\hbar\omega_x\left(n_x+\frac12\right)
			+\hbar\omega_y\left(n_y+\frac12\right).
		\]
		把 $q$ 当作连续参数。定态的 Ehrenfest 方程给出
		\[
			m\omega_x^2\langle x\rangle=q\mathcal E_x,\qquad
			m\omega_y^2\langle y\rangle=q\mathcal E_y.
		\]
		Hellmann--Feynman 定理于是给出
		\begin{align*}
			\frac{\partial E}{\partial q}
			&=\left\langle\frac{\partial H}{\partial q}\right\rangle
			=-\mathcal E_x\langle x\rangle-\mathcal E_y\langle y\rangle\\
			&=-\frac{q}{m}\left(\frac{\mathcal E_x^2}{\omega_x^2}
			+\frac{\mathcal E_y^2}{\omega_y^2}\right).
		\end{align*}
		从 $0$ 积分到物理电荷 $q$，得
		\[
			\boxed{E_{n_xn_y}(q)=E_{n_xn_y}(0)
			-\frac{q^2}{2m}\left(\frac{\mathcal E_x^2}{\omega_x^2}
			+\frac{\mathcal E_y^2}{\omega_y^2}\right)}.
		\]
	\end{enumerate}
}

	
\item 	
(共30分)设哈密顿算符的矩阵表示为$H=H_{0}+\lambda H^{\prime}$，
$H_{0}=\left(\begin{array}{cc}E_{1} & 0 \\ 0 & E_{2}\end{array}\right)$， $H^{\prime}=\left(\begin{array}{cc}0 & i a \\ -i a & 0\end{array}\right)$，
其中$E_{1}<E_{2}$，$\left|\frac{\lambda a}{E_{2}-E_{1}}\right|<<1$。
\begin{enumerate}
	%\renewcommand{\labelenumi}{\arabic{enumi}.}
	% A(\Alph) a(\alph) I(\Roman) i(\roman) 1(\arabic)
	%设定全局标号series=example	%引用全局变量resume=example
	%[topsep=-0.3em,parsep=-0.3em,itemsep=-0.3em,partopsep=-0.3em]
	%可使用leftmargin调整列表环境左边的空白长度 [leftmargin=0em]
	\item
以$ \lambda H ^{\prime}  $作为微扰,用微扰论求$ H $的准确至二级的本征值以及准确至一级的归一化本征矢。
	
	
\item 
精确求解$ H $的本征值,并与近似结果比较。


	
	
	
	
\end{enumerate}


\banswer{
	令 $\Delta=E_2-E_1>0$、$\epsilon=\lambda a/\Delta$，并以
	$|1\rangle=(1,0)^T$、$|2\rangle=(0,1)^T$ 为无扰基。
	\begin{enumerate}
		\item 一级能量修正均为零，二级修正为
		\[
			E_1^{(2)}=-\frac{\lambda^2a^2}{\Delta},\qquad
			E_2^{(2)}=\frac{\lambda^2a^2}{\Delta}.
		\]
		所以
		\[
			\boxed{\mathcal E_1=E_1-\frac{\lambda^2a^2}{\Delta}
			+O(\lambda^4)},\qquad
			\boxed{\mathcal E_2=E_2+\frac{\lambda^2a^2}{\Delta}
			+O(\lambda^4)}.
		\]
		准确到一级、并用一个不影响该精度的因子归一化，本征矢可写为
		\[
			\boxed{|\mathcal E_1\rangle=
			\frac{1}{\sqrt{1+\epsilon^2}}
			\begin{pmatrix}1\\i\epsilon\end{pmatrix}},\qquad
			\boxed{|\mathcal E_2\rangle=
			\frac{1}{\sqrt{1+\epsilon^2}}
			\begin{pmatrix}i\epsilon\\1\end{pmatrix}}.
		\]

		\item 精确本征值为
		\[
			\boxed{\mathcal E_\pm=\frac{E_1+E_2}{2}
			\pm\sqrt{\frac{\Delta^2}{4}+\lambda^2a^2}}.
		\]
		把平方根按 $|\lambda a/\Delta|\ll1$ 展开，负号支与正号支分别给出上一问的
		$E_1-\lambda^2a^2/\Delta$ 和
		$E_2+\lambda^2a^2/\Delta$。
	\end{enumerate}
}



\item 
(共30分)电子$ 1 $和$ 2 $分别局域于两个空间格点上。设体系哈密顿量为
$$
H=-2 C\left(s_{x}^{(1)} s_{x}^{(2)}+s_{y}^{(1)} s_{y}^{(2)}\right)
$$
其中, $\vec{s}^{(i)} (i=1,2)$为自旋算符,$ C $是一个常数。令$ \hbar=1 $。

\begin{enumerate}
	%\renewcommand{\labelenumi}{\arabic{enumi}.}
	% A(\Alph) a(\alph) I(\Roman) i(\roman) 1(\arabic)
	%设定全局标号series=example	%引用全局变量resume=example
	%[topsep=-0.3em,parsep=-0.3em,itemsep=-0.3em,partopsep=-0.3em]
	%可使用leftmargin调整列表环境左边的空白长度 [leftmargin=0em]
	\item
写出体系的总自旋$\vec{s}=\vec{s}^{(1)}+\vec{s}^{(2)}$的平方算符$\vec{s}^{2}$和$S_{z}=S_{z}^{(1)}+S_{z}^{(2)}$的可取值。

\item 
用$\vec{s}^{2}$和$ s_{z} $表示体系的哈密顿量。

\item 
求出体系所有能级的能量。

\item 
沿$ z $方向施加一磁场,强度为$ B $,求出体系所有能级的能量。


\end{enumerate}


\banswer{
	这里 $\hbar=1$，每个电子满足
	$(\vec s^{(i)})^2=3/4$。
	\begin{enumerate}
		\item 两个自旋 $1/2$ 可耦合成
		\[
			s=1,\quad s^2=s(s+1)=2,\quad m=s_z=1,0,-1,
		\]
		或
		\[
			s=0,\quad s^2=0,\quad m=0.
		\]

		\item 利用
		$s_{1z}s_{2z}=[s_z^2-s_{1z}^2-s_{2z}^2]/2$，可得
		\[
			s_{1x}s_{2x}+s_{1y}s_{2y}
			=\frac12(s^2-s_z^2-1).
		\]
		故
		\[
			\boxed{H=C(1+s_z^2-s^2)}.
		\]

		\item 所有能级为
		\[
			\boxed{E_{1,\pm1}=0\quad(\text{二重})},\qquad
			\boxed{E_{1,0}=-C},\qquad
			\boxed{E_{0,0}=C}.
		\]

		\item 若 Zeeman 项写成 $H_B=-\gamma B s_z$，总能量为
		\[
			E_{s,m}=C[1+m^2-s(s+1)]-\gamma Bm.
		\]
		即
		\[
			\boxed{E_{1,1}=-\gamma B,\quad E_{1,0}=-C,\quad
			E_{1,-1}=\gamma B,\quad E_{0,0}=C}.
		\]
		若题目采用相反的磁矩符号，只需令 $\gamma\to-\gamma$。
	\end{enumerate}
}

	
\end{enumerate}
